% Inbuilt themes in beamer

\documentclass[aspectratio=169]{beamer}
\usepackage[export]{adjustbox}
\usepackage{xcolor}

\definecolor{darkred}{RGB}{176, 31, 36}
\definecolor{darkgreen}{RGB}{44, 95, 47}

\usepackage{listings}
\usepackage{booktabs}
\usepackage{graphicx}

% Title page details: 
\title{Tema 3.2: Redes recurrentes}
\subtitle{GRU, LSTM y Bidireccionales}

% Theme choice:
\usetheme{Madrid}
\usecolortheme{seagull}
\useinnertheme{circles}

\author{Pablo Pérez Núñez}
 
\setbeamertemplate{navigation symbols}{}
\setbeamertemplate{author}{}

% Ajusta las dimensiones de la diapositiva

\setbeamertemplate{blocks}[rounded][shadow=false]
% \setbeamertemplate{blocks}[default]
\definecolor{myblue}{RGB}{27, 122, 194}
\definecolor{lightblue}{RGB}{55, 150, 225}
\definecolor{lightbluetwo}{RGB}{164, 209, 242}

\setbeamercolor{block title}{bg=lightblue,fg=white}
\setbeamercolor{block body}{bg=lightbluetwo,fg=black}

\setbeamercolor{title}{fg=white, bg =lightblue}
\setbeamercolor{frametitle}{fg=white, bg = lightblue}

\setbeamercolor*{palette secondary}{fg=white,bg=lightblue}

\defbeamertemplate*{footline}{infolines}
{
  \leavevmode%
  \hbox{%
  \begin{beamercolorbox}[wd=\paperwidth,ht=2.25ex,dp=1ex,left]{title in head/foot}%
    \hspace*{2ex} \usebeamerfont{title in head/foot}\insertshorttitle \ - \insertsubtitle
  \end{beamercolorbox}
  % \begin{beamercolorbox}[wd=.25\paperwidth,ht=2.25ex,dp=1ex,right]{date in head/foot}%
  %   \usebeamerfont{date in head/foot}\hspace*{2em}
  %   \insertframenumber{} / \inserttotalframenumber\hspace*{2ex} 
  %\end{beamercolorbox}
  }%
  \vskip0pt%
}

\setbeamertemplate{title page}{
  \centering
  % \usefonttheme{serif}
  \begin{minipage}{0.95\textwidth} % Ajustar el ancho según sea necesario

    \vspace{1em}
    {\color{myblue}\rule{\linewidth}{2pt}\par}
    \vspace{.5cm}
    {\bfseries \huge \textcolor{myblue}{\inserttitle}\par }
    \vspace{1em}
    {\bfseries \LARGE \textcolor{myblue}{\insertsubtitle}\par }
      
    \vspace{2.5em}

    \begin{columns}
      \column{0.6\textwidth}
      \vspace{6em}\\
      {\bfseries \large \textcolor{myblue}{Aprendizaje profundo / Deep Learning}} \\
      \vspace{1em}
      {\textcolor{myblue}{\insertauthor}}
      \column{0.4\textwidth}
      \vfill
      \centering
      \includegraphics[height=4.5em]{imgs/LogoAEPIA.pdf} \\
      \vspace{1em}
      \includegraphics[height=2.5em]{imgs/UIMP.png}
    \end{columns}

    \vspace{1em}

    {\color{myblue}\rule{\linewidth}{.5pt}\par}

  \end{minipage}
    

  %  \usefonttheme{default}
}

% Define the template for section slides
\AtBeginSection[]{
  \begin{frame}
    \vfill
    \LARGE\insertsectionhead\par
    {\color{myblue}\rule{\linewidth}{1pt}\par}
    \vfill
  \end{frame}
}
 

% MIS COSAS --------------------------------
% Bloque con margen inferior
\newenvironment{blockm}[1]{%
  \begin{block}{\textbf{#1}}%
  }{%
  \end{block}%
  \vspace{1em}%
}
% Dos imágenes
\newcommand{\twoimages}[3]{%
  \begin{columns}
    \begin{column}{0.5\textwidth}
      \centering
      \includegraphics[width=#3\textwidth]{#1}
    \end{column}
    \begin{column}{0.5\textwidth}
      \centering
      \includegraphics[width=#3\textwidth]{#2}
    \end{column}
  \end{columns}
}

% Comando para mostrar el título de la subsección al inicio de cada nueva subsección
\AtBeginSubsection[]
{
  \begin{frame}
    \vfill
    \small\insertsectionhead\par
    {\color{myblue}\rule{\linewidth}{1pt}\par}
	\LARGE\insertsubsectionhead\par
    \vfill
  \end{frame}
}

% ------------------------------------------

\begin{document}

% Title page frame
\begin{frame}[plain]
	\titlepage 
\end{frame}

\logo{}


\subsection{Gate Recurrent Unit (GRU)}

\begin{frame}{Gate Recurrent Unit (GRU)}
	\begin{blockm}{Definición}
		Redes recurrentes con `gates' (puertas) que regulan el flujo de información que pasa de una etapa a la siguiente.
	\end{blockm}
	\begin{itemize}
		\item \textbf{Regulación:} Transmitir desde cada neurona una salida ponderada entre su activación y la entrada recurrente.
		\begin{itemize}
			\item El factor de ponderación se aprende y suele ser en muchos casos próximo a 0 o a 1.
		\end{itemize}
		\item \textbf{Efecto memoria:} Puede ocurrir que entre la etapa $j$ y la etapa $k$ fluya la información de la etapa anterior, es decir, 
				$\mathbf{h}_{\mathbf{i}}^{(j)} \simeq \mathbf{h}_{\mathbf{i}}^{(j+1)} \simeq \ldots \simeq \mathbf{h}_{\mathbf{i}}^{(k-2)} \simeq \mathbf{h}_{\mathbf{i}}^{(k-1)}$.	Esto provoca que la salida en la etapa $k$ dependa en gran medida de la etapa $j$, que puede ser muy anterior.
	\end{itemize}
	Intuitivamente, una red GRU puede permitir que las neuronas se especialicen en determinadas propiedades de la secuencia que están tratando para mejorar la predicción.
\end{frame}

\begin{frame}{RNN vs GRU}

	\begin{columns}
		\begin{column}{0.5\textwidth}
			\includegraphics[width=.9\textwidth, center]{imgs/tema4/rnn/RNNsimple.pdf}
		\end{column}

		\begin{column}{0.5\textwidth}
			\includegraphics[width=.85\textwidth, center]{imgs/tema4/rnn/GRU.pdf}
		\end{column}
	\end{columns}
\end{frame}

\begin{frame}{GRU con relevancia}
	\begin{columns}
		\begin{column}{0.5\textwidth}
			\includegraphics[width=.85\textwidth, center]{imgs/tema4/rnn/GRU.pdf}
		\end{column}

		\begin{column}{0.5\textwidth}
			\includegraphics[width=.85\textwidth, center]{imgs/tema4/rnn/GRUrele.pdf}
		\end{column}
	\end{columns}
\end{frame}

\begin{frame}{Notación}
	\begin{columns}
		\begin{column}{0.4\textwidth}
			\includegraphics[width=\textwidth, center]{imgs/tema4/rnn/GRUrele.pdf}
		\end{column}

		\begin{column}{0.6\textwidth}
			\begin{equation*}
				\begin{aligned}
				& \boldsymbol{\Gamma}_{\mathrm{u}} = \sigma(\mathbf{M}_{u}[\mathbf{h}^{(t-1), \mathbf{x}^{(t)}}]+\mathbf{b}_{u})\\
				& \boldsymbol{\Gamma}_{\mathrm{r}} = \sigma(\mathbf{M}_{r}[\mathbf{h}^{(t-1), \mathbf{x}^{(t)}}]+\mathbf{b}_{r})\\
				& \boldsymbol{\hat{h}}^{(t)} = \text{tanh}(\mathbf{M}[\boldsymbol{\Gamma}_{\mathrm{r}}\odot\mathbf{h}^{(t-1)}, \mathbf{x}^{(t)}]+\mathbf{b}_{h}) \\
				& \boldsymbol{h}^{(t)} = \boldsymbol{\Gamma}_{\mathrm{r}}\odot\mathbf{\hat{h}}^{(t)}+(1-\boldsymbol{\Gamma}_{\mathrm{u}}) \odot \mathbf{h}^{(t-1)} \\[5px]
				& \text{donde:} \\[5px]
				& \mathbf{M} = [\mathbf{W}, \mathbf{U}] \\
				& \mathbf{M}_{r} = [\mathbf{W}_{r}, \mathbf{U}_{r}] \\
				& \mathbf{M}_{u} = [\mathbf{W}_{u}, \mathbf{U}_{u}] \\
				\end{aligned}
			\end{equation*}
		\end{column}
	\end{columns}
\end{frame}


\begin{frame}{Notación}
	\begin{columns}
		\begin{column}{0.4\textwidth}
			\includegraphics[width=\textwidth, center]{imgs/tema4/rnn/GRUrele2.pdf}
		\end{column}

		\begin{column}{0.6\textwidth}
			\includegraphics[width=.9\textwidth, center]{imgs/tema4/rnn/GRUrele2desp.pdf}
		\end{column}
	\end{columns}
\end{frame}

\subsection{Long Short Term Memory (LSTM)}


\begin{frame}{Long Short Term Memory (LSTM)}

	\begin{blockm}{Descripción}
		Se presentaron antes que las GRU (Hochreiter y Schmidhuber, 1997). Más complejas y flexibles que las redes GRU ya que cuentan con varias `gates'.
	\end{blockm}
	
	\begin{columns}
		\begin{column}{0.4\textwidth}
			En las redes LSTM hay dos vectores de información que se pasan de una etapa a otra:

			\begin{itemize}
				\item $ \mathbf{h}^{(t)} $: Salida de la capa LSTM.
				\item $ \mathbf{c}^{(t)} $: \textit{Estado} de la capa LSTM.
			\end{itemize}
		\end{column}

		\begin{column}{0.6\textwidth}
			\includegraphics[width=\textwidth, center]{imgs/tema4/rnn/GRUsec.pdf}
		\end{column}
	\end{columns}
\end{frame}

\begin{frame}{Long Short Term Memory (LSTM)}

	\begin{blockm}{Descripción}
		Se presentaron antes que las GRU (Hochreiter y Schmidhuber, 1997). Más complejas y flexibles que las redes GRU ya que cuentan con varias `gates'.
	\end{blockm}
	
	\begin{columns}
		\begin{column}{0.4\textwidth}
			En las redes LSTM hay dos vectores de información que se pasan de una etapa a otra:

			\begin{itemize}
				\item $ \mathbf{h}^{(t)} $: Salida de la capa LSTM.
				\item $ \mathbf{c}^{(t)} $: \textit{Estado} de la capa LSTM.
			\end{itemize}
		\end{column}

		\begin{column}{0.6\textwidth}
			\includegraphics[width=\textwidth, center]{imgs/tema4/rnn/LSTMsec.pdf}
		\end{column}
	\end{columns}
\end{frame}

\begin{frame}{Notación}
	\begin{columns}
		\begin{column}{0.6\textwidth}
			\includegraphics[width=\textwidth, center]{imgs/tema4/rnn/LSTM.pdf}
		\end{column}

		\begin{column}{0.4\textwidth}
			\begin{equation*}
				\begin{aligned}
				& \boldsymbol{\Gamma}_{\mathrm{u}} = \sigma(\mathbf{M}_{u}[\mathbf{h}^{(t-1)}, \mathbf{x}^{(t)}]+\mathbf{b}_{u})\\
				& \boldsymbol{\Gamma}_{\mathrm{f}} = \sigma(\mathbf{M}_{f}[\mathbf{h}^{(t-1)}, \mathbf{x}^{(t)}]+\mathbf{b}_{f})\\
				& \boldsymbol{\Gamma}_{\mathrm{o}} = \sigma(\mathbf{M}_{o}[\mathbf{h}^{(t-1)}, \mathbf{x}^{(t)}]+\mathbf{b}_{o})\\
				& \boldsymbol{\hat{c}}^{(t)} = \text{tanh}(\mathbf{M}_{c}[\mathbf{h}^{(t-1)}, \mathbf{x}^{(t)}]+\mathbf{b}_{c}) \\
				& \boldsymbol{c}^{(t)} = \boldsymbol{\Gamma}_{\mathrm{u}}\odot\boldsymbol{\hat{c}}^{(t)} + \boldsymbol{\Gamma}_{\mathrm{f}} \odot \boldsymbol{c}^{(t-1)} \\
				& \boldsymbol{h}^{(t)} = \boldsymbol{\Gamma}_{\mathrm{o}}\odot\text{tanh}(\boldsymbol{c}^{(t)}) \\[5px]
				& \text{donde:} \\[5px]
				& \mathbf{M}_{u} = [\mathbf{W}_{u}, \mathbf{U}_{u}] \hspace{1em} \mathbf{M}_{f} = [\mathbf{W}_{f}, \mathbf{U}_{f}]\\
				& \mathbf{M}_{o} = [\mathbf{W}_{o}, \mathbf{U}_{o}] \hspace{1em} \mathbf{M}_{c} = [\mathbf{W}_{c}, \mathbf{U}_{c}]\\
				\end{aligned}
			\end{equation*}
		\end{column}
	\end{columns}
\end{frame}


\subsection{Bidireccionales}

\begin{frame}{RNN Bidireccionales}

	\begin{blockm}{Descripción}
		Las redes bidireccionales (Schuster y Paliwal, 1997) tienen en cuenta el contexto \textbf{posterior} a un elemento en una secuencia.
	\end{blockm}
	
	El contexto puede ser relevante:
	\begin{itemize}
		\item \textit{``Deberías ver {\color{blue}Roma}, {\color{teal}es una ciudad} con gran cantidad de monumentos''}
		\item \textit{``Deberías ver {\color{red}Roma}, {\color{teal}es una película} de Alfonso Cuarón que te va a encantar''}
	\end{itemize}
	En este caso es necesario utilizar elementos posteriores a la palabra Roma para determinar de que estamos hablando.

\end{frame}

\begin{frame}{RNN Bidireccional (BRNN)}
	Constituidas por:
	\begin{itemize}
		\item \textbf{Capa forward}: Celdas RNN que procesan la secuencia de inicio a fin.
		\item \textbf{Capa backward}: Celdas que la procesan en orden inverso.
	\end{itemize}
	\vspace{1em}
	\includegraphics[width=\textwidth, center]{imgs/tema4/rnn/RNNBidi mini.pdf}
\end{frame}

\begin{frame}{RNN Bidireccional (BRNN)}
	Capas forward y backward en detalle:\\
	\vspace{1em}
	\includegraphics[width=.9\textwidth, center]{imgs/tema4/rnn/RNNBidi.pdf}
\end{frame}

\begin{frame}[t]{Referencias}
	\begin{enumerate}
	  \item \textbf{Redes recurrentes (GRU, LSTM y Bidireccionales), Oscar Luaces}
	\end{enumerate}
  \end{frame}

\end{document}
